\part{Gaudin Model}

\section{Commuting Hamiltonians}

\begin{exercise}[The rational Gaudin spin system and its invariant interaction]
Let $V_{s_i}$ carry spin $s_i\in\frac12\Z_{\ge0}$, with
$[S_i^a,S_i^b]=\ii\varepsilon_{abc}S_i^c$,
$\sum_a(S_i^a)^2=s_i(s_i+1)\id$, and operators at different
sites commuting. These are the $q\to1$ representations of the
section Representation theory. For pairwise distinct real
parameters $z_1,\dots,z_N$ define
\[
 \Omega_{ij}=\mathbf S_i\cdot\mathbf S_j,\qquad
 G_i=\sum_{j\ne i}\frac{\Omega_{ij}}{z_i-z_j}
 \quad\text{on }\mathcal H=\bigotimes_iV_{s_i}.
\]
Intrinsically, $\Omega_{ij}$ is the contraction of the two spin
actions with the invariant Euclidean pairing on the physical
rotation Lie algebra. A Gaudin Hamiltonian is
$H=\sum_i c_iG_i$ for specified real $c_i$.
\begin{enumerate}[label=(\alph*)]
\item Prove $[\Omega_{ij},\Omega_{k\ell}]=0$ for disjoint
pairs and $[\Omega_{ij},\Omega_{ik}+\Omega_{jk}]=0$ for
distinct $i,j,k$, using invariance of the pairing.
Combine these identities with partial fractions to prove
$[G_i,G_j]=0$.
\item Prove $\sum_iG_i=0$ and
\[
 \sum_i z_iG_i=\frac12\left(
 \mathbf S_{\rm tot}^2-\sum_i s_i(s_i+1)\id\right).
\]
Show that every $G_i$ is self-adjoint and rotation invariant.
\item Use joint spectral projections to express the partition
function of $H$ as a sum over joint energies with their
multiplicities.
\end{enumerate}
\end{exercise}

\section{Bethe equations}

\begin{exercise}[Gaudin Bethe states and cancellation of unwanted terms]
Let $\Omega$ be a tensor product of highest-weight states and put
$S^-(w)=\sum_iS_i^-/(w-z_i)$,
$S_i^\pm=S_i^x\pm\ii S_i^y$.
For distinct roots $w_1,\dots,w_M$ avoiding all $z_i$ define
$\Psi=\prod_aS^-(w_a)\Omega$.
\begin{enumerate}[label=(\alph*)]
\item Compute the action for $M=0,1$ directly from the spin
commutation relations. Then commute $G_i$ successively through
the product to prove that, if $\Psi\ne0$ and
\[
 \sum_i\frac{s_i}{w_a-z_i}
       -\sum_{b\ne a}\frac1{w_a-w_b}=0
       \quad(1\le a\le M),
\]
it has joint eigenvalues
\[
 g_i=\sum_{j\ne i}\frac{s_i s_j}{z_i-z_j}
           -\sum_a\frac{s_i}{z_i-w_a}.
\]
Keep the terms in which one $S^-(w_a)$ is replaced by $S_i^-$;
their coefficients are the displayed Bethe equations.
\item Compute $S^+_{\rm tot}\Psi$ by the same method and show
that it vanishes under these equations. Its total spin is
$s_{\rm tot}=\sum_i s_i-M$ when the vector is nonzero.
Check both identities of the section Commuting Hamiltonians
against the resulting eigenvalues.
\item For two spins $s_1=s_2=\tfrac12$ solve the $M=1$
equation, show $w=(z_1+z_2)/2$, and identify the alternating
spin state without choosing a decomposition into eigenvectors.
\end{enumerate}
\end{exercise}

\section{Opers}

\begin{exercise}[Projective differential operators and the Miura transformation]
On a Riemann surface let $K$ be the holomorphic cotangent line.
A projective second-order differential operator acts from
$K^{-1/2}$ to $K^{3/2}$, with principal symbol one and
vanishing first-derivative term in a local coordinate:
$D=\partial_z^2-t(z)$. On the sphere use its fixed spin line
$K^{1/2}=\mathcal O(-1)$, where $\mathcal O(-1)$ denotes the
tautological line of the section Cohomology and $K$-theory.
Such an operator, or its projective
solution local system, is the scalar form of a $PGL_2$-oper.
Meromorphic poles are permitted at specified marked points.
For the Gaudin parameters define the Miura function
\[
 \chi(z)=\sum_i\frac{s_i}{z-z_i}-\sum_a\frac1{z-w_a},
 \qquad
 D=(\partial_z-\chi)(\partial_z+\chi)
   =\partial_z^2-(\chi^2-\chi').
\]
\begin{enumerate}[label=(\alph*)]
\item Use $T_L\mathbb P(W)=\Hom(L,W/L)$ for $\dim W=2$
and a fixed identification $\bigwedge^2W=\C$ to prove that
the cotangent line is the square of the tautological line.
Under $z=z(\zeta)$ transform a solution as a
half-density and derive the transformation rule
$t_\zeta=(z')^2t_z-\frac12\{z,\zeta\}$, where
$\{z,\zeta\}=z'''/z'-\frac32(z''/z')^2$.
Verify compatibility on overlapping coordinates.
\item Expand $\chi^2-\chi'$ at each $w_a$. Show that its
double pole cancels automatically and its simple pole
vanishes exactly when the Bethe equation holds. At a marked
point prove
\[
 t(z)=\frac{s_i(s_i+1)}{(z-z_i)^2}
                  +\frac{2g_i}{z-z_i}+O(1).
\]
Recover the joint Hamiltonian eigenvalues from these coefficients.
\item Set
$y(z)=\prod_a(z-w_a)\prod_i(z-z_i)^{-s_i}$ on a simply
connected domain avoiding the marked points. Verify $Dy=0$
and construct a second solution $y\int y^{-2}\dd z$.
Use the Bethe equations to show that $y^{-2}\dd z$ has
zero residues at the $w_a$; it is regular at each $z_i$.
Deduce projectively trivial monodromy for the constructed
oper when $2s_i$ are nonnegative integers.
\end{enumerate}
\end{exercise}

\section{Langlands duality}

\begin{exercise}[Rank-one root data and the dual group of the spin symmetry]
For a complex algebraic torus $T$, its character and cocharacter
lattices are $X^*(T)=\Hom(T,\C^\times)$ and
$X_*(T)=\Hom(\C^\times,T)$, paired by composition.
A root datum consists of these paired lattices and paired finite
sets of roots and coroots with $\langle\alpha,\alpha^\vee\rangle=2$,
such that the reflections
$x\mapsto x-\langle x,\alpha^\vee\rangle\alpha$ and their
dual reflections preserve the respective root sets.
Its dual interchanges the two lattices and the two root sets.
\begin{enumerate}[label=(\alph*)]
\item Use the specified maximal torus
$T=\{\operatorname{diag}(t,t^{-1})\}\subset SL_2$.
Compute $X^*(T)=\Z$, root $2$, $X_*(T)=\Z$, and coroot $1$.
After interchange identify the result with the torus and root
data of $PGL_2=GL_2/\C^\times$.
\item The indicial exponents of the oper at $z_i$ solve
$r(r-1)=s_i(s_i+1)$. Compute them as $-s_i$ and $s_i+1$
and their difference as $2s_i+1$. Explain this integer
through the highest weight $2s_i$ and the rank-one
half-sum of positive roots, with the positive root fixed above.
\item Show that changing signs of a lift of the oper's
two-dimensional solution transport has no effect on its
projectivization. Relate this to the quotient from $SL_2$
to the dual group $PGL_2$.
Compare the constructed Bethe-to-oper map with \cite{frenkel-opers}.
\end{enumerate}
\end{exercise}

\section{Geometric Langlands}

\begin{exercise}[Line bundles, Hecke modifications, and the abelian eigencondition on the sphere]
For a smooth complex curve $C$, a rank-one local system is a
functor from its fundamental groupoid to one-dimensional complex
vector spaces. A divisor is a finite integer sum of points.
For a divisor $D$ define $\mathcal O(D)$ by its meromorphic
sections $f$ satisfying $\operatorname{ord}_x(f)+D(x)\ge0$
wherever the section is required to be holomorphic.
The divisor Picard group is divisors modulo the divisors
of nonzero meromorphic functions; use its line bundles
$\mathcal O(D)$ as the Picard space in this exercise.
A Hecke modification at $x$ is
$\mathcal O(D)\mapsto\mathcal O(D+x)$.
For a local system $E$ on $C$, the abelian Hecke eigencondition
for a family of lines $\mathcal A$ on $\operatorname{Pic}(C)$ is
\[
 h^*\mathcal A\simeq E\boxtimes\mathcal A,\qquad
 h(x,L)=L(x),
\]
with the isomorphisms compatible with successive modifications.
Use the coarse Picard space and $C=\mathbb P^1$ in this exercise.
\begin{enumerate}[label=(\alph*)]
\item On the spectral sphere with fixed affine parameter $z$
prove $\operatorname{div}(z-a)=(a)-(\infty)$.
Deduce that every divisor is equivalent to $d(\infty)$,
where $d$ is its degree, and that degree is unchanged by
a principal divisor. Prove that multiplication by a
meromorphic function realizes the corresponding isomorphism
of divisor line bundles. Thus compute the divisor Picard
space as $\{\mathcal O(d(\infty)):d\in\Z\}$.
\item Use simple connectivity of the sphere to identify every
rank-one local system with a constant line $E_0$.
Set $\mathcal A_d=E_0^{\otimes d}$, using dual lines for
negative $d$, and construct the Hecke eigenisomorphism.
Prove compatibility for modifications at two different
points from the tensor symmetry.
\item Fixing $\mathcal A_0=\C$, prove that the eigencondition
determines all $\mathcal A_d$ and the compatible identifications.
Identify the dual torus as $\C^\times$ by interchanging its
character and cocharacter lattices. State explicitly the
curve, group, and rank for this abelian correspondence.
\end{enumerate}
\end{exercise}

\begin{exercise}[Prediction: two-site Gaudin spectroscopy from oper residues]
Consider two spin-$\tfrac12$ sites with distinct real $z_1,z_2$
and physical Hamiltonian
$H_{\rm phys}=J(z_1-z_2)G_1=J\,\mathbf S_1\cdot\mathbf S_2$,
$J>0$. The measured energy transfer is $\hbar\omega$.
\begin{enumerate}[label=(\alph*)]
\item For $M=0$ calculate the oper and its residue at $z_1$.
For $M=1$ insert $w=(z_1+z_2)/2$ and calculate the changed
residue. Use $E=J(z_1-z_2)g_1$ to derive the energies
$J/4$ and $-3J/4$.
\item Use total spin to determine the multiplicities three
and one. Derive a singlet-triplet absorption line at
$\hbar\omega=J$ and the partition function
$Z=e^{3\beta J/4}+3e^{-\beta J/4}$.
Verify agreement with the two-site quiver prediction after
its stated additive energy shift.
\item Compute the response to
$S_1^z-S_2^z$: it maps the normalized singlet into the
zero-magnetization triplet with norm one. Determine the
zero-temperature spectral weight of the absorption line.
State the isolated-pair, isotropic-exchange, and negligible
line-broadening assumptions. Distinguish the derived
two-site Bethe-to-oper calculation from the separate
abelian Hecke eigencondition on the sphere.
\end{enumerate}
\end{exercise}
